%%
%% fall-lecture11.tex
%% 
%% Made by Alex Nelson <pqnelson@gmail.com>
%% Login   <alex@lisp>
%% 
%% Started on  2026-10-24T12:04:41-0700
%% Last update 2026-10-24T12:04:41-0700
%% 

\lecture{}

\begin{theorem}[Pontryagin]
Let $B$ be a compact $(n+1)$-dimensional smooth manifold with boundary
$M=\boundary B$. Then the Stiefel--Whitney numbers of $M$ are all
zero. (Note: neither $B$ nor $M$ are required to be connected.)
\end{theorem}

\begin{proof}[Sketchy proof sketch]
It amounts to looking at the long exact sequence for cohomology, then
using exactness to get this vanishing.
\end{proof}

\begin{theorem}[Thom]
Let $M$ be a closed $n$-dimensional smooth manifold suchthat all the
Stiefel--Whitney numbers of $M$ are zero. Then $M$ is the boundary of
some compact smooth $(n+1)$-manifold $B$: $M=\boundary B$.
\end{theorem}

\begin{remark}
Thom's theorem requires more machinery and is beyond the scope of this
course, but its proof can be found in the
textbook~\cite[Th~4.10]{milnor1974characteristic}. %60
\end{remark}

\begin{definition}
Let $M_{1}$ and $M_{2}$ be closed $n$-dimensional smooth manifolds.
We say $M_{1}$ and $M_{2}$ \define{belong to the Cobordism Class} if
there exists a smooth $(n+1)$-manifold $W$ whose boundary is their
disjoint union $M_{1}\sqcup M_{2}=\boundary W$. 
\end{definition}

\begin{remark}
(1) Note we do not require $M_{1}$, $M_{2}$, or $W$ to be connected.

(2) We do not require $M_{1}$, $M_{2}$, or $W$ to be orientable.
\end{remark}

\begin{remark}
There is a stronger notion of an \define{Oriented Cobordism Class}: 
we say the \emph{orientable} smooth $n$-dimensional manifolds
$M_{1}$ and $M_{2}$ belong to the same oriented cobordism class
if there exists an \emph{orientable} smooth $(n+1)$-manifold $W$ such
that $\boundary W = M_{1}\sqcup(-M_{2})$.
\end{remark}

\begin{xca}
Prove ``\dots belong to the same unoriented cobordism class'' is an
equivalence relation. Is the same true for ``\dots belong to the same
oriented cobordism class''?
\end{xca}

\begin{corollary}
Let $M_{1}$, $M_{2}$ be smooth $n$-dimensional manifolds.
Then $M_{1}$ and $M_{2}$ belong to the same unoriented cobordism
class if and only if all the Stiefel--Whitney numbers of $M_{1}$ and
$M_{2}$ agree.
\end{corollary}
% TODO
(I'm omitting the proof, it is in my notebook on page 106--107.)

\begin{example}
Let $M_{1}=\sphere{2}$, let $M_{2}=\torus{2}$. Notice
$w(\T\sphere{2})=1$. What is $w(\T M_{2})=w_{0}+w_{1}+w_{2}$? We know
$w_{0}=1$. Since $\torus{2}$ is orientable, $w_{1}=0$ and we know
$\chi(\torus{2})=0$ so $w_{2}=0$. This means $w(\T M_{2})=1$.

There are multiple possible $W$ such that $\boundary W=M_{1}\sqcup M_{2}$.
For example, just take the disjoint union of the ball $B^{3}$ and the
solid torus $\sphere{1}\times\disk{2}$. Then $\boundary(B^{3}\sqcup\sphere{1}\times\disk{2})=M_{1}\sqcup M_{2}$.
This is boring, it's disconnected. What's a \emph{connected} manifold
$W$ such that $\boundary W=M_{1}\sqcup M_{2}$?

Well, we could take the solid torus $\sphere{1}\times\disk{2}$, then
surgically remove a suitably small open ball $B^{3}\subset\sphere{1}\times\disk{2}$.
This would work.
\end{example}

\begin{question}
Can you find examples $M_{1}$ and $M_{2}$ such that if such a $W$
(satisfying $\boundary W=M_{1}\sqcup M_{2}$) exists, then $W$
\emph{must} be connected? (The Stiefel--Whitney numbers for $M_{1}$,
$M_{2}$ may not all be zero, unlike the preceding example.)
\end{question}